\begin{theorem}[S-37: Split semihomogeneous replacements retain a trace kernel for every finite action]
\label{thm:semireg-rm-split-semihomogeneous}
Let $(X,\theta,F)$ be a real-multiplication abelian
fourfold of type $(2,2)$, with the balanced class $v_q$
of S-22 and $q>0$. Let
\[
 E=\bigoplus_{j=1}^{s}Q_j[n_j]
 \tag{S37.1}
\]
be a finite direct sum of shifts of simple vector bundles
such that
\[
 \operatorname{ch}(Q_j)=r_j e^{\ell_j},\qquad
 r_j=\operatorname{rk}(Q_j)>0,\qquad
 \ell_j=\frac{c_1(Q_j)}{r_j}\in F\theta.
 \tag{S37.2}
\]
In particular, this includes simple semihomogeneous
bundles with slopes in $F\theta$. Suppose $E$ has
positive virtual rank $N$ and
$e^B\operatorname{ch}(E)=Nv_q$ for a rational algebraic
divisor class $B$.
For any lift $\widetilde E$ under any effective finite
coherent translation--tensor action $G$ as in S-21,
\[
 \operatorname{ext}^2_{\operatorname{Perf}(X)^G}
                  (\widetilde E,\widetilde E)\ge24,
 \qquad
 \dim\ker\sigma^2_{\widetilde E}\ge3.
 \tag{S37.3}
\]
Thus no such lift is fully semiregular. Neither
simplicity nor gluability of $\widetilde E$ is assumed.
The assertion concerns the direct sums (S37.1); it
does not assert persistence through arbitrary nonsplit
extensions or cones of these bundles.
\end{theorem}
\begin{proof}
The semihomogeneous specialization of (S37.2) is Mukai,
\emph{Abelian variety and spin representation}, equation
(4.8), p.~10, in the author's sixteen-page version.
The argument below only needs the explicit character
condition (S37.2).

Put $w_j=(-1)^{n_j}r_j$. The degree-two character
identity gives
$B=-N^{-1}\sum_jw_j\ell_j\in F\theta$.
Write the normalized slopes as
\[
 B+\ell_j=x_j\theta_1+y_j\theta_2,
 \qquad x_j,y_j\in\mathbf R.
\]
Since $\theta_1^3=\theta_2^3=0$, comparing the
coefficients of $\theta_1^a\theta_2^b/(a!b!)$
for $0\le a,b\le2$ gives the moment identity
\[
 \sum_j w_j
 \begin{pmatrix}1\\x_j\\y_j\\x_jy_j\end{pmatrix}
 \begin{pmatrix}1&x_j&y_j&x_jy_j\end{pmatrix}
 =\operatorname{diag}(N,-Nq,-Nq,Nq^2).
 \tag{S37.4}
\]
All nine monomials used here are linearly independent
in the exterior algebra of the two eigenspace factors.
The matrix on the right has rank four and signature
$(2,2)$. Aggregate the weights belonging to the same
pair $(x_j,y_j)$. At least four distinct pairs have
nonzero aggregate weight, since each pair contributes
a matrix of rank at most one. In fact at least two
aggregate weights are positive and two negative, by
the signature, although only the rank is needed below.

We must retain these slope blocks for an arbitrary
equivariant linearization, which may mix summands.
Let $A=\operatorname{End}_X(E)$ and $J=\operatorname{rad}(A)$.
The shifted simple bundles are indecomposable objects
with endomorphism ring $\mathbf C$. Krull--Schmidt
decomposition therefore identifies $A/J$ with a product
of matrix algebras indexed by their isomorphism classes;
see Krause, \emph{Krull--Schmidt categories and projective
covers}, Theorem~4.2 and Corollary~4.4, pp.~9--10,
including the radical description immediately following.
For each distinct normalized slope $s=(x,y)$, let
$\bar e_s$ be the central idempotent selecting all
classes with that slope, allowing all shifts.
Translations and tensoring with degree-zero line bundles
preserve the slope and preserve the shift of a nonzero
bundle. Consequently the induced $G$-action fixes every
$\bar e_s$; this holds for every linearization of $E$.

Taking invariants is exact in characteristic zero, so
\[
 A^G\longrightarrow(A/J)^G
 \tag{S37.5}
\]
is surjective, with nilpotent kernel $J\cap A^G$.
The family of orthogonal idempotents $\bar e_s$ can
therefore be lifted to orthogonal idempotents
$e_s\in A^G$ summing to one. The exact noncommutative
algebra input is Aguiar--Mahajan, \emph{Topics in
Hyperplane Arrangements}, Lemma~D.29, p.~533;
Lemma~D.27 on the same page identifies any such lift
with the original slope summand. These lemmas apply
to nilpotent ideals without commutativity assumptions.
Thus we obtain an equivariant direct sum
\[
 \widetilde E=\bigoplus_s\widetilde E_s,
 \qquad
 \operatorname{rk}(E_s)=
      w_s:=\sum_{j:(x_j,y_j)=s}w_j.
 \tag{S37.6}
\]
For the rank assertion, the lifted idempotent and the
original slope projector have the same image in $A/J$.
The isomorphic-idempotent conclusion of Lemma~D.27
therefore makes their images isomorphic as underlying
perfect objects. No equivariant choice of individual
simple-bundle summands was required.

For any slope with $w_s\ne0$, the natural map
\[
 H^2(X,\mathcal O_X)\longrightarrow
       \operatorname{Ext}^2_X(E_s,E_s),
 \qquad \eta\longmapsto\eta\operatorname{id}_{E_s}
 \tag{S37.7}
\]
is injective. Composing with the ordinary perfect trace
gives multiplication by the nonzero virtual rank $w_s$.
Its image is $G$-invariant: translations act trivially
on $H^2(\mathcal O_X)$, and tensor operations and the
linearizations commute with this natural cohomology
action. Hence (S37.7) injects into the equivariant
Ext group. Since $h^2(\mathcal O_X)=6$, each of the
at least four nonzero-rank blocks supplies six
independent degree-two classes. The diagonal Ext
groups of the direct sum (S37.6) are direct summands
of its full self-Ext group, proving the first inequality
in (S37.3).

Let $m$ be that full equivariant Ext dimension and
let $b=\dim\ker\sigma^2_{\widetilde E}$.
The invariant CY4 Serre pairing and the full
trace--Atiyah transpose identity of S-21 give
\[
 \operatorname{rank}(\sigma^2_{\widetilde E})
   =\operatorname{rank}(\operatorname{ob}^2_{\widetilde E})
   =m-b.
\]
Their composition has rank eighteen by the balanced
character calculation in S-22 and the identity-sector
Hochschild comparison of S-21. These rank identities
do not require vanishing of negative self-Ext groups.
The rank inequality for a composition therefore gives
\[
 18\ge2(m-b)-m=m-2b.
 \tag{S37.8}
\]
As $m\ge24$, this yields $b\ge3$ and proves the theorem.
\end{proof}
